\section{Lamport clocks and vector clocks}\label{sec:lc}
Handout: LAB1 §3.4 \\
Code: \TODO{\src{protocols/logical-clocks} or \src{protocols/lamport}}

% ---------------------------------------------------------------------
\subsection{Algorithm}
Every node $i$ keeps a scalar $L$ and a vector $V$ of length $n$.
\begin{center}\small
\begin{tabular}{@{}lll@{}}
\toprule
Event at node $i$ & Lamport & Vector \\ \midrule
local or send & $L \leftarrow L+1$ & $V[i] \leftarrow V[i]+1$ \\
receive $(L_m, V_m)$ & $L \leftarrow \max(L, L_m)+1$ &
  $V[k] \leftarrow \max(V[k], V_m[k])\ \forall k$, then $V[i] \leftarrow V[i]+1$ \\
\bottomrule
\end{tabular}
\end{center}
Total-order timestamp: \TODO{$L.i$, or $n\cdot L + i$ -- write the one your code uses}.

Properties to verify:
\begin{itemize}[itemsep=1pt]
  \item $e \to e' \Rightarrow L(e) < L(e')$, but not the converse.
  \item $V(e) < V(e') \iff e \to e'$; otherwise the events are concurrent ($e \parallel e'$).
  \item The total order is arbitrary for concurrent events.
\end{itemize}

% ---------------------------------------------------------------------
\subsection{Implementation}
\begin{itemize}[itemsep=2pt]
  \item Clock update rules: \TODO{function names + links}
  \item Message piggybacking (what is sent with each message): \TODO{}
  \item \fld{Compare}: \TODO{where, and what it returns}
  \item Concurrency: \TODO{which mutex protects the clock; see the variant in
        Section~\ref{sec:lc-variant}}
\end{itemize}

\phead{State fields and inputs}
\begin{table}[H]\centering\small
\begin{tabularx}{\linewidth}{@{}lX@{}}
\toprule
Field / input & Meaning \\ \midrule
\TODO{\fld{lamport}}  & \TODO{} \\
\TODO{\fld{vector}}   & \TODO{} \\
\TODO{\fld{events}}   & \TODO{} \\
\TODO{\fld{Local}}    & \TODO{} \\
\TODO{\fld{Send}}     & \TODO{} \\
\TODO{\fld{Compare}}  & \TODO{} \\
\TODO{\fld{Auto}}     & \TODO{} \\
\bottomrule
\end{tabularx}
\end{table}

% ---------------------------------------------------------------------
\subsection{Scenario coverage}
\begin{table}[H]\centering\small
\begin{tabularx}{\linewidth}{@{}llXl@{}}
\toprule
Report section & Handout & What it shows & Export \\ \midrule
\ref{sec:lc-slide}  & A         & lecture slide 37: $L(b) > L(e)$ but $b \parallel e$ & \TODO{file} \\
\ref{sec:lc-equal}  & B         & three concurrent events with equal $L$ & \TODO{file} \\
\ref{sec:lc-auto}   & C         & invariant over random events (\fld{Auto}) & \TODO{file} \\
\ref{sec:lc-crash}  & (fault A) & node crash and recovery & \TODO{file} \\
\ref{sec:lc-part}   & (fault B) & blocked link: send without receive & \TODO{file} \\
\ref{sec:lc-variant}& variant   & \TODO{} & \TODO{file} \\
\bottomrule
\end{tabularx}
\end{table}

% ---------------------------------------------------------------------
\scenario{Normal (handout A): the example on lecture slide 37}{lc-slide}
\phead{Setup}
\begin{setuptable}
Nodes / latency & 3 nodes, \TODO{}\,ms, Docker \\
Export          & \TODO{} \\
\end{setuptable}
\phead{Steps}
\begin{steps}
  \item \TODO{the exact sequence of Local / Send inputs that reproduces the figure}
\end{steps}

\begin{figure}[H]\centering
\begin{tikzpicture}[x=1cm,y=0.9cm]
  \proc{3}{node-1}{10}
  \proc{1.5}{node-2}{10}
  \proc{0}{node-3}{10}
  % TODO: adjust positions/labels to match your run
  \node[ev,label=above:{\scriptsize a}] (a) at (1,3) {};
  \node[ev,label=above:{\scriptsize b}] (b) at (3,3) {};
  \node[ev,label=above:{\scriptsize c}] (c) at (5,1.5) {};
  \node[ev,label=above:{\scriptsize d}] (d) at (6.5,1.5) {};
  \node[ev,label=below:{\scriptsize e}] (e) at (2,0) {};
  \node[ev,label=below:{\scriptsize f}] (f) at (8.5,0) {};
  \draw[msg] (b) -- node[lbl,above right]{m1} (c);
  \draw[msg] (d) -- node[lbl,above right]{m2} (f);
\end{tikzpicture}
\caption{Expected event structure (\TODO{check against slide 37}).}
\end{figure}

\phead{Result}
\begin{table}[H]\centering\small
\begin{tabular}{@{}llllrl@{}}
\toprule
Event & Node & Kind & $L$ & $V$ & total order \\ \midrule
a & node-1 & local   & \fillin & \fillin & \fillin \\
b & node-1 & send    & \fillin & \fillin & \fillin \\
c & node-2 & receive & \fillin & \fillin & \fillin \\
d & node-2 & send    & \fillin & \fillin & \fillin \\
e & node-3 & local   & \fillin & \fillin & \fillin \\
f & node-3 & receive & \fillin & \fillin & \fillin \\
\bottomrule
\end{tabular}
\end{table}
\phead{\fld{Compare} results}
\begin{table}[H]\centering\small
\begin{tabular}{@{}llll@{}}
\toprule
Pair & Lamport says & Vector says & Agree? \\ \midrule
$(b, e)$ & \fillin & \fillin & \fillin \\
$(b, f)$ & \fillin & \fillin & \fillin \\
\TODO{} & & & \\
\bottomrule
\end{tabular}
\end{table}
\screenshot{lc-slide37.png}{Space-time diagram reproducing slide 37.}
\phead{Explanation}
\TODO{3--5 sentences. Why $L(b) > L(e)$ does not mean $e \to b$.}

% ---------------------------------------------------------------------
\scenario{Handout B: concurrent events with equal Lamport values}{lc-equal}
\phead{Steps}
\begin{steps}
  \item On a fresh run, have node-1, node-2, node-3 each produce one local event (no messages).
  \item \fld{Compare} each pair.
\end{steps}
\phead{Result}
\TODO{table: three events, $L=1$, vectors, total-order numbers, pairwise comparison}
\phead{Explanation}
\TODO{3--5 sentences. Does the order 1.1 $<$ 1.2 $<$ 1.3 have any physical meaning?}

% ---------------------------------------------------------------------
\scenario{Handout C: invariant after random events}{lc-auto}
\phead{Setup}
\begin{setuptable}
Nodes / latency & \TODO{} \\
Auto settings   & \TODO{rate, duration} \\
Export          & \TODO{} \\
\end{setuptable}
\phead{Result}
\begin{table}[H]\centering\small
\begin{tabular}{@{}lllll@{}}
\toprule
Event 1 & Event 2 & Lamport order & Vector relation & \fld{agrees} \\ \midrule
\fillin & \fillin & \fillin & \fillin & \fillin \\
\fillin & \fillin & \fillin & \fillin & \fillin \\
\fillin & \fillin & \fillin & \fillin & \fillin \\
\bottomrule
\end{tabular}
\end{table}
\phead{Explanation}
\TODO{3--5 sentences. Every causal pair agrees; some Lamport-ordered pairs are concurrent.}

% ---------------------------------------------------------------------
\scenario{Fault A: a node crashes and recovers}{lc-crash}
\phead{Steps}
\begin{steps}
  \item \TODO{produce events, crash node-2, send to it, recover, produce more events}
\end{steps}
\phead{Explanation}
\TODO{3--5 sentences. Are the counters restored (\fld{sdk.Save}/\fld{sdk.Load})?
What would go wrong if a recovered node restarted at $L=0$?}

% ---------------------------------------------------------------------
\scenario{Fault B: blocked link or high latency}{lc-part}
\phead{Steps}
\begin{steps}
  \item \TODO{block node-1$\to$node-2, Send, heal, Send again}
\end{steps}
\phead{Explanation}
\TODO{3--5 sentences. A send event exists without a receive event; the vector
clocks still describe causality correctly.}

% ---------------------------------------------------------------------
\scenario{Design variant: \TODO{name}}{lc-variant}
\hint{The handout gives no fixed variant here. Options: (1) on receive, skip the
$\max$ and just increment -- then $e \to e'$ no longer implies $L(e) < L(e')$;
(2) merge vectors without $\max$ (copy the incoming vector) -- causality is lost;
(3) remove the mutex around the clock (your \nolinkurl{protocols/lamport} code has
this experiment commented in) and show the corrupted state under concurrent sends.}
\phead{Code change}
\TODO{}
\phead{Result}
\TODO{baseline vs.\ variant}
\phead{Explanation}
\TODO{3--5 sentences.}

% ---------------------------------------------------------------------
\subsection{Discussion questions}
\begin{qa}{Why does $L(e) < L(e')$ not imply $e \to e'$?}
\TODO{2--4 sentences.}
\end{qa}
\begin{qa}{Does the total order among concurrent events have physical meaning?}
\TODO{2--4 sentences.}
\end{qa}
\begin{qa}{What does a vector clock cost compared with a Lamport clock?}
\TODO{2--4 sentences (message size $O(n)$, fixed membership).}
\end{qa}

\subsection{Boundaries (what is not implemented)}
\begin{itemize}[itemsep=1pt]
  \item \TODO{e.g.\ fixed membership, no causal-delivery buffer, no total-order broadcast.}
\end{itemize}
